% Generated by tools/build_new_dists_anthology_sections.py; do not edit by hand.

\section{Association-scheme permutation groups}\label{proof:new-dist:association_schemes}

\resultbox{\textbf{Canonical testing artifacts.}\par\smallskip Exact backend: \path{new_dists/association_scheme_sigma_mgfs/verification.py}.\par Regression: \path{new_dists/association_scheme_sigma_mgfs/test_verification.py}.\par Marimo: \path{marimo_notebooks/association_schemes.py}.}

\subsection*{Scope and outcome}

For every finite automorphism group $G\leq\Sym(\Omega)$, we prove that the
complete permutation $\sigma$-MGF on $\C[\Omega]$ is determined by the
power-fixed-point data $F_r(g)=|\Fix_\Omega(g^r)|$.  We then prove explicit
fixed-point templates for homogeneous spaces, affine translation actions, and
wreath-product actions.  These specialize to Johnson, Hamming, Grassmann,
form, polar, building, graph, Doob, folded-cube, and halved-cube families.
Stable ranges and degree-two obstructions are kept separate from exact finite
formulas.  An accompanying marimo notebook constructs representative matrix
groups and compares direct matrix traces, cycle products, and literal
configuration orbits in 99 exact coefficient checks.  Every check passes.

\subsection{The universal power-fixed-point theorem}

Let a finite group $G$ act on a finite set $\Omega$, let $V=\C[\Omega]$, and
let $P_g$ denote the resulting permutation matrix.  For a composition or
partition $\tau=(\tau_1,\ldots,\tau_s)$ put
\[
 W_\tau(V)=\bigotimes_{j=1}^s\Sym^{\tau_j}V,
 \qquad a_\tau(G,\Omega)=\dim W_\tau(V)^G.
\]
Our convention is
\[
 \MGF_{G,\Omega}(\mathbf t)
 =\E_{g\in G}\ExpS(X_V(g)h_1)
 =\sum_\tau a_\tau(G,\Omega)m_\tau(\mathbf t).
\]
For $g\in G$ define
\[
 F_r(g)=|\Fix_\Omega(g^r)|,
 \qquad c_d(g)=\#\{\text{$d$-cycles of $g$ on $\Omega$}\}.
\]

\begin{theorem}[Association-scheme power-fixed-point theorem]
For every finite permutation action,
\begin{align}
 F_r(g)&=\sum_{d\mid r}d\,c_d(g),\label{nd:association_schemes:eq:fixed-from-cycles}\\
 c_d(g)&=\frac1d\sum_{r\mid d}\mu(d/r)F_r(g),\label{nd:association_schemes:eq:mobius}\\
 \MGF_{G,\Omega}(\mathbf t)
 &=\frac1{|G|}\sum_{g\in G}
   \prod_{i\geq1}\prod_{d\geq1}(1-t_i^d)^{-c_d(g)},\label{nd:association_schemes:eq:universal}\\
 [m_\tau]\MGF_{G,\Omega}
 &=\#\left(G\backslash\prod_j\Sym^{\tau_j}\Omega\right).
 \label{nd:association_schemes:eq:orbits}
\end{align}
Thus the complete $\sigma$-MGF is determined by the collection of numbers
$F_r(g)$.
\end{theorem}

\begin{proof}
A point on a $d$-cycle is fixed by $g^r$ exactly when $d\mid r$, proving
\eqref{nd:association_schemes:eq:fixed-from-cycles}; arithmetic M\"obius inversion gives
\eqref{nd:association_schemes:eq:mobius}.  The eigenvalues of a $d$-cycle permutation matrix are all
$d$th roots of unity, so
\[
 \det(I-tP_g\mid\C^d)=1-t^d.
\]
Multiplication over cycles and variables, followed by Reynolds averaging,
proves \eqref{nd:association_schemes:eq:universal}.  Finally, the standard monomial basis of
$\Sym^a\C[\Omega]$ is indexed by $a$-element multisets in $\Omega$.  Hence
$W_\tau(V)$ is a permutation module on the product in \eqref{nd:association_schemes:eq:orbits}; its
orbit sums form a basis of invariants.
\end{proof}

\begin{remark}[What is tested directly]
The verification stores $P_g$ exactly by the location of the single nonzero
entry in each column.  It computes symmetric-power characters from the actual
matrix traces using Newton's recurrence
\[
 d\,h_d(P_g)=\sum_{r=1}^d\operatorname{Tr}(P_g^r)h_{d-r}(P_g),\qquad h_0=1.
\]
This matrix route, the determinant cycle-factor route, and literal
configuration orbits are separate computations of the target coefficient.
\end{remark}

\subsection{Three reusable fixed-point templates}

\subsubsection{Homogeneous and parabolic actions}

Let $\Omega=G/H$ with left multiplication.

\begin{proposition}[Coset fixed points]\label{nd:association_schemes:prop:coset}
For $x\in G$,
\begin{equation}\label{nd:association_schemes:eq:coset}
 \boxed{|\Fix_{G/H}(x)|
 =\frac{|C_G(x)|\,|x^G\cap H|}{|H|}.}
\end{equation}
\end{proposition}

\begin{proof}
The coset $yH$ is fixed precisely when $y^{-1}xy\in H$.  For each element of
$x^G\cap H$, exactly $|C_G(x)|$ choices of $y$ conjugate $x$ to it.  Dividing
the number of such $y$ by $|H|$ accounts for the representatives of a coset.
Apply this to $x=g^r$, then use \eqref{nd:association_schemes:eq:mobius}--\eqref{nd:association_schemes:eq:universal}.
\end{proof}

This covers Grassmannians and other parabolic quotients, dual-polar spaces,
generalized polygons from $BN$-pairs, and building residues.

\subsubsection{Affine translation actions}

Let $A$ be a finite additive group, $H\leq\Aut(A)$, and
$G=A\rtimes H$ act by $x\mapsto b+h(x)$.  Define
\[
 b_r=(I+h+\cdots+h^{r-1})b.
\]

\begin{proposition}[Affine kernel/image formula]\label{nd:association_schemes:prop:affine}
For $g=(b,h)$,
\begin{equation}\label{nd:association_schemes:eq:affine}
 \boxed{F_r(b,h)=
 \begin{cases}
 |\ker(I-h^r)|,&b_r\in\operatorname{im}(I-h^r),\\
 0,&b_r\notin\operatorname{im}(I-h^r).
 \end{cases}}
\end{equation}
\end{proposition}

\begin{proof}
Induction gives $g^r(x)=b_r+h^r(x)$.  Thus $x$ is fixed exactly when
$(I-h^r)x=b_r$.  A fiber of the homomorphism $I-h^r$ is empty outside its
image and is a coset of its kernel otherwise.
\end{proof}

\subsubsection{Product and wreath actions}

Let $H\curvearrowright\Gamma$ and let $H\wr S_d=H^d\rtimes S_d$ act on
$\Gamma^d$.  Write $g=(h_1,\ldots,h_d;\pi)$.  For a cycle
$C=(i_1\ldots i_\ell)$ of $\pi$, let
\[
 a_C=h_{i_\ell}\cdots h_{i_1},\qquad e_C(r)=\gcd(\ell,r),
 \qquad f_s(a)=|\Fix_\Gamma(a^s)|.
\]

\begin{proposition}[Marked-cycle product formula]\label{nd:association_schemes:prop:wreath}
\begin{equation}\label{nd:association_schemes:eq:wreath}
 \boxed{F_r(g)=\prod_{C\in\operatorname{Cyc}(\pi)}
 f_{r/e_C(r)}(a_C)^{e_C(r)}.}
\end{equation}
\end{proposition}

\begin{proof}
Under $\pi^r$, an $\ell$-cycle splits into $e=\gcd(\ell,r)$ coordinate
cycles.  Transport once around any one has monodromy conjugate to
$a_C^{r/e}$.  Each of the $e$ cycles admits $f_{r/e}(a_C)$ choices.  Different
cycles of $\pi$ use disjoint coordinates, so their counts multiply.
\end{proof}

\subsection{Johnson schemes}

Let $\Omega=\binom{[v]}k$ and $G=S_v$.  Suppose $\sigma$ has cycle type
$\lambda=1^{m_1}2^{m_2}\cdots$.  A $j$-cycle of $\sigma$ splits under
$\sigma^r$ into $e=\gcd(j,r)$ cycles of length $j/e$.  A fixed $k$-subset is
a union of these cycles, hence
\begin{equation}\label{nd:association_schemes:eq:johnson-fixed}
 \boxed{F_r^{J(v,k)}(\lambda)=[u^k]\prod_{j\geq1}
 \left(1+u^{j/\gcd(j,r)}\right)^{m_j\gcd(j,r)}.}
\end{equation}
With $c_d^J(\lambda)=d^{-1}\sum_{r\mid d}\mu(d/r)F_r^J(\lambda)$ and
$z_\lambda=\prod_jj^{m_j}m_j!$,
\begin{equation}\label{nd:association_schemes:eq:johnson-smg}
 \boxed{\MGF_{J(v,k)}=\sum_{\lambda\vdash v}\frac1{z_\lambda}
 \prod_{i,d\geq1}(1-t_i^d)^{-c_d^J(\lambda)}.}
\end{equation}

For $D=|\tau|$, a configuration contributing to $[m_\tau]$ uses at most
$kD$ underlying points.  Relabelling those points proves
\begin{equation}\label{nd:association_schemes:eq:johnson-stable}
 [m_\tau]\MGF_{J(v,k)}\text{ is constant for }v\geq kD.
\end{equation}
The stable objects are isomorphism classes of $\tau$-colored $k$-uniform
multihypergraphs without isolated vertices.  Ordered pairs are classified by
intersection size, so
\begin{equation}\label{nd:association_schemes:eq:johnson-pair}
 \boxed{[m_{(1,1)}]\MGF_{J(v,k)}=\min(k,v-k)+1.}
\end{equation}
Thus no coefficientwise limit exists when $\min(k,v-k)\to\infty$.

\subsection{Hamming and Doob product schemes}

For $\Omega=[q]^d$ and $G=S_q\wr S_d$, Proposition~\ref{nd:association_schemes:prop:wreath} gives
\begin{equation}\label{nd:association_schemes:eq:hamming-fixed}
 F_r(g)=\prod_C\left|\Fix_{[q]}
 \left(a_C^{r/e_C(r)}\right)\right|^{e_C(r)}.
\end{equation}
Consequently
\begin{equation}\label{nd:association_schemes:eq:hamming-smg}
 \boxed{\MGF_{H(d,q)}=\frac1{(q!)^dd!}\sum_{\pi\in S_d}
 \sum_{h_1,\ldots,h_d\in S_q}\prod_{i,\ell\geq1}
 (1-t_i^\ell)^{-\frac1\ell\sum_{r\mid\ell}\mu(\ell/r)F_r(g)}.}
\end{equation}

For fixed $d$ and $D=|\tau|$, once $q\geq D$ the orbit of the $D$ symbols in
each coordinate is exactly its equality partition in the partition lattice
$\Pi_D$.  If $K_\tau=\prod_jS_{\tau_j}$ permutes repeated objects within the
colored multiset blocks, then
\begin{equation}\label{nd:association_schemes:eq:hamming-stable}
 \boxed{[m_\tau]\MGF_{H(d,q)}
 =\#\bigl((K_\tau\times S_d)\backslash\Pi_D^d\bigr),\qquad q\geq D.}
\end{equation}
On the other hand, ordered word pairs are classified by Hamming distance:
\begin{equation}\label{nd:association_schemes:eq:hamming-pair}
 \boxed{[m_{(1,1)}]\MGF_{H(d,q)}=d+1.}
\end{equation}

For the Doob graph $D(m,n)=\mathrm{Sh}^{\square m}\square K_4^{\square n}$,
take
\[
 G_{m,n}=(\Aut(\mathrm{Sh})\wr S_m)\times(S_4\wr S_n).
\]
Fixed points multiply between the two Cartesian blocks.  The Shrikhande
vertex action has four ordered-pair types and $S_4\curvearrowright[4]$ has
two.  Permuting equal factors leaves multisets of types, so
\begin{equation}\label{nd:association_schemes:eq:doob-pair}
 \boxed{[m_{(1,1)}]\MGF_{D(m,n)}=\binom{m+3}{3}(n+1).}
\end{equation}
In contrast, $H(2m+n,4)$ has pair coefficient $2m+n+1$ even though the two
families have the same distance-regular parameters.

\subsection{Grassmann, polar, polygon, and building actions}

\subsubsection{Grassmann schemes}

Let $\Omega=\Gr_k(\F_q^n)$ and take the projective image of $\GL_n(q)$.  Directly,
\begin{equation}\label{nd:association_schemes:eq:grass-fixed}
 F_r(g)=\#\{U\leq\F_q^n:\dim U=k,\ g^rU=U\}.
\end{equation}
Equivalently, with $P=P_{k,n-k}$, Proposition~\ref{nd:association_schemes:prop:coset} gives
\[
 F_r(g)=\frac{|C_G(g^r)|\,|(g^r)^G\cap P|}{|P|}.
\]
Invariant-subspace counts may be computed from the rational canonical form
of $g^r$.  Ordered pairs are classified by $\dim(U\cap W)$:
\begin{equation}\label{nd:association_schemes:eq:grass-pair}
 \boxed{[m_{(1,1)}]\MGF_{\Gr_k(n,q)}=\min(k,n-k)+1.}
\end{equation}
The span of $D$ participating $k$-spaces has dimension at most $kD$, and an
isomorphism between supported configurations extends to the ambient space.
Thus
\begin{equation}\label{nd:association_schemes:eq:grass-stable}
 [m_\tau]\MGF_{\Gr_k(n,q)}\text{ stabilizes for }n\geq kD
 \quad(q,k,\tau\text{ fixed}).
\end{equation}
Fixed-rank field-size limits can fail: four points on a projective line carry
a cross-ratio parameter.

\subsubsection{Dual-polar spaces and generalized polygons}

Let $V$ be a finite formed space of Witt rank $d$ and let $\Omega=G/P$ be the
maximal isotropic or singular subspaces.  Then
\begin{equation}\label{nd:association_schemes:eq:dualpolar-fixed}
 F_r(g)=\frac{|C_G(g^r)|\,|(g^r)^G\cap P|}{|P|}.
\end{equation}
Pair orbits are indexed by intersection codimension, hence
\begin{equation}\label{nd:association_schemes:eq:dualpolar-pair}
 \boxed{[m_{(1,1)}]\MGF_{\mathrm{DP}(d,q)}=d+1.}
\end{equation}

The same coset formula applies to a flag-transitive generalized polygon with
point set $G/P$.  If the point action is distance-transitive, its pair orbits
are indexed by distance.  Therefore the pair coefficients for generalized
quadrangles, hexagons, and octagons are respectively $3,4,5$.  Without
flag/distance transitivity, the order parameters alone do not imply these
values and do not determine the full series.

\subsubsection{Buildings}

Suppose $G$ has a $BN$-pair with Weyl group $W$.  For chambers $G/B$, Bruhat
decomposition gives $B\backslash G/B\cong W$, so
\begin{equation}\label{nd:association_schemes:eq:building-chamber}
 \boxed{[m_{(1,1)}]\MGF_{G,G/B}=|W|.}
\end{equation}
For residues of type $J$,
\begin{equation}\label{nd:association_schemes:eq:building-residue}
 \boxed{[m_{(1,1)}]\MGF_{G,G/P_J}=|W_J\backslash W/W_J|.}
\end{equation}
These conclusions require the natural strong or Weyl transitivity.  Higher
coefficients need not be controlled by $W$ because triples of flags can carry
field-valued moduli.

\subsection{Translation schemes of forms}

All families in this section use Proposition~\ref{nd:association_schemes:prop:affine}.  The target
representation is the complex permutation module on the additive form space;
the finite-field matrices are used to construct its zero-one permutation
matrices.

\subsubsection{Bilinear forms}

Let $A=M_{m\times n}(\F_q)$ and let $(P,Q)$ act linearly by
$L(X)=PXQ^{-1}$.  For $g=(B,P,Q)$ define
$B_r=(I+L+\cdots+L^{r-1})B$.  Then
\begin{equation}\label{nd:association_schemes:eq:bilinear-fixed}
 \boxed{F_r(g)=
 \begin{cases}
 q^{\dim\ker(I-L^r)},&B_r\in\operatorname{im}(I-L^r),\\
 0,&\text{otherwise}.
 \end{cases}}
\end{equation}
After translating one member of a pair to zero, differences are classified by
rank:
\begin{equation}\label{nd:association_schemes:eq:bilinear-pair}
 \boxed{[m_{(1,1)}]\MGF_{\mathrm{Bil}(m,n,q)}=\min(m,n)+1.}
\end{equation}
For fixed $m,q,\tau$, translating one of $D=|\tau|$ matrices to zero leaves
combined row support of dimension at most $m(D-1)$.  Thus the coefficient is
stable once $n\geq m(D-1)$.

\subsubsection{Alternating, Hermitian, and quadratic forms}

For $A=\operatorname{Alt}_n(\F_q)$ use $L(X)=PXP^T$; for Hermitian matrices
over $\F_{q^2}$ use $L(X)=PXP^{\sigma T}$.  Formula \eqref{nd:association_schemes:eq:bilinear-fixed}
applies verbatim with the appropriate linear operator.  Rank classifies the
one-form orbits in these two cases, while alternating ranks are even.  Hence
\begin{align}
 [m_{(1,1)}]\MGF_{\mathrm{Alt}(n,q)}&=\left\lfloor\frac n2\right\rfloor+1,
 \label{nd:association_schemes:eq:alt-pair}\\
 [m_{(1,1)}]\MGF_{\mathrm{Herm}(n,q)}&=n+1.
 \label{nd:association_schemes:eq:herm-pair}
\end{align}

Let $A=\mathcal Q(n,q)$ be the additive space of quadratic forms, with
$P\cdot Q=Q\circ P^{-1}$.  Again \eqref{nd:association_schemes:eq:affine} is exact.  For odd $q$,
each positive rank has two congruence types and rank zero has one, yielding
\begin{equation}\label{nd:association_schemes:eq:quadratic-pair}
 \boxed{[m_{(1,1)}]\MGF_{\mathcal Q(n,q)}=2n+1\qquad(q\text{ odd}).}
\end{equation}
Even characteristic requires the usual rank/type refinement but no change to
the affine fixed-point formula.  The displayed pair coefficients already rule
out rank-growing coefficientwise limits.  Field-size limits can also fail in
higher degree: after translating a triple of bilinear forms and making one
difference invertible, the other retains similarity invariants such as
eigenvalues.

\subsection{Affine-polar, folded-cube, and halved-cube actions}

For an affine-polar action take the additive formed space $A=V$ and the
relevant isometry, similitude, or semisimilitude group $H$.  Equation
\eqref{nd:association_schemes:eq:affine} is the complete exact answer.  For $D=|\tau|$, translate one
of the $D$ points to zero.  The other $D-1$ differences span a formed subspace
of dimension at most $D-1$.  Its restricted form and labelled vectors
determine the orbit; Witt extension embeds equivalent supported configurations
into all sufficiently large ambient spaces.  Therefore
\begin{equation}\label{nd:association_schemes:eq:affpolar-stable}
 [m_\tau]\MGF_{\mathrm{AffPol}(N,q)}\text{ is eventually constant as }N\to\infty
 \quad(q,\tau\text{ fixed}).
\end{equation}

For the folded cube, let
\[
 A_n=\F_2^n/\langle\mathbf1\rangle,\qquad G=A_n\rtimes S_n.
\]
For $g=(b,\sigma)$, equation \eqref{nd:association_schemes:eq:affine} becomes
\begin{equation}\label{nd:association_schemes:eq:folded-fixed}
 F_r(g)=\begin{cases}
 2^{\dim\ker(I-\sigma^r\mid A_n)},&b_r\in\operatorname{im}(I-\sigma^r),\\
 0,&\text{otherwise}.
 \end{cases}
\end{equation}
Weights $w$ and $n-w$ represent the same difference orbit, so
\begin{equation}\label{nd:association_schemes:eq:folded-pair}
 \boxed{[m_{(1,1)}]\MGF_{\mathrm{FoldCube}_n}=\left\lfloor\frac n2\right\rfloor+1.}
\end{equation}

For the halved cube use
$E_n=\{x\in\F_2^n:\sum_i x_i=0\}$ and $G=E_n\rtimes S_n$.
Formula \eqref{nd:association_schemes:eq:folded-fixed} holds with $A_n$ replaced by $E_n$.
Difference weights are the even integers from zero to the largest even value
at most $n$, giving
\begin{equation}\label{nd:association_schemes:eq:halved-pair}
 \boxed{[m_{(1,1)}]\MGF_{\frac12Q_n}=\left\lfloor\frac n2\right\rfloor+1.}
\end{equation}

\subsection{Distance-regular and strongly regular graphs}

For any finite graph $\Gamma$, equations
\eqref{nd:association_schemes:eq:mobius}--\eqref{nd:association_schemes:eq:universal} apply with
$G=\Aut(\Gamma)$ and $\Omega=V(\Gamma)$.  If the action is
distance-transitive of diameter $D$, then
\begin{equation}\label{nd:association_schemes:eq:distance-pair}
 [m_{(1,1)}]\MGF_\Gamma=D+1.
\end{equation}
More generally, the pair coefficient is the permutation rank of the vertex
action.

The $4\times4$ rook graph and the Shrikhande graph both have strongly regular
parameters $(16,6,2,2)$, but their full vertex actions have ranks three and
four.  In the Cayley model on $\mathbb Z_4^2$, the Shrikhande vertex
stabilizer has orbits
\[
 \{0\},\quad \{\text{six neighbors}\},\quad
 \{(2,0),(0,2),(2,2)\},\quad \{\text{remaining six vertices}\}.
\]
Therefore the $\sigma$-MGF is not determined by the intersection array,
spectrum, or strongly regular parameters.  At the other extreme, a graph with
trivial automorphism group on $N$ vertices has
\[
 \MGF_\Gamma=\prod_{i\geq1}(1-t_i)^{-N}.
\]

\subsection{Asymptotic classification}

Let $D=|\tau|$.  The proof mechanism behind every positive result is bounded
support: a fixed-degree configuration lives in a bounded subobject and ambient
automorphisms extend isomorphisms of that subobject.  Every negative result in
the table is already forced by an unbounded degree-two coefficient.

\begin{center}
\small
\renewcommand{\arraystretch}{1.12}
\begin{tabularx}{\linewidth}{@{}>{\raggedright\arraybackslash}p{0.47\linewidth}X@{}}
\toprule
Family and regime & Coefficientwise behavior \\
\midrule
$J(v,k)$, fixed $k$, $v\to\infty$ & Stable for $v\geq kD$ \\
$J(v,k)$, $\min(k,v-k)\to\infty$ & No limit; pair coefficient diverges \\
$H(d,q)$, fixed $d$, $q\to\infty$ & Stable for $q\geq D$ \\
$H(d,q)$, fixed $q$, $d\to\infty$ & No limit; pair coefficient $d+1$ \\
$\Gr_k(n,q)$, fixed $q,k$, $n\to\infty$ & Stable for $n\geq kD$ \\
Grassmann, fixed rank, $q\to\infty$ & Can fail through cross-ratio moduli \\
Bilinear, fixed $m,q$, $n\to\infty$ & Stable for $n\geq m(D-1)$ \\
Alternating/Hermitian/quadratic, rank $\to\infty$ & No limit in degree two \\
Dual-polar, Witt rank $\to\infty$ & No limit; pair coefficient $d+1$ \\
Affine-polar, fixed $q$, rank $\to\infty$ & Eventually stable for fixed $\tau$ \\
Building chambers, rank $\to\infty$ & No limit when $|W|\to\infty$ \\
Bounded residues in classical towers, fixed $q$ & Stable by bounded support \\
Doob/folded/halved dimension $\to\infty$ & No limit in degree two \\
\bottomrule
\end{tabularx}
\end{center}

\subsection{Exact computational verification}

\subsubsection{Construction and comparison strategy}

The companion module \texttt{verification.py} uses only exact finite
arithmetic for the structural checks.
\begin{enumerate}
\item Enumerate the stated finite matrix group (or a faithful permutation
image) and the finite set $\Omega$.
\item Construct $P_g$ by applying each group element to every point of
$\Omega$.  The dimension is exactly $|\Omega|$.
\item Compute $F_r(g)$ directly from $P_g^r$ and independently from the
appropriate coset, affine, invariant-subspace, or wreath formula.
\item Recover $c_d(g)$ by M\"obius inversion and compare with literal cycle
decomposition.\par
\item Compute $[m_\tau]\MGF$ from Newton matrix traces, from cycle-factor
coefficient extraction, and from direct orbits on
$\prod_j\Sym^{\tau_j}\Omega$.
\end{enumerate}

The marimo notebook presents these checks by group family.  Existing Johnson,
Hamming, Grassmann, and polar constructors are imported from the canonical
association-scheme package in this workspace; the new notebook adds the
affine-form, building, graph, Doob, folded, and halved constructors.  No
combined anthology PDF is modified.

\subsubsection{Representative exact cases}

\begin{center}
\small
\renewcommand{\arraystretch}{1.1}
\begin{longtable}{@{}p{0.38\linewidth}rrrrl@{}}
\toprule
Action & $|\Omega|$ & $|H|$ & $|G|$ & $[m_{11}]$ & Audit \\
\midrule
\endfirsthead
\toprule
Action & $|\Omega|$ & $|H|$ & $|G|$ & $[m_{11}]$ & Audit \\
\midrule
\endhead
$S_v$ on $k$-subsets: $(4,2),(5,2),(5,1)$ & $6,10,5$ & -- & $48,120,120$ & $3,3,2$ & exact \\
$S_q\wr S_d$: $(2,2),(2,3),(3,2)$ & $4,8,9$ & -- & $8,48,72$ & $3,4,3$ & exact \\
$PGL_3(2)$ on $\Gr_1,\Gr_2$; $PGL_2(3)$ on $\Gr_1$ & $7,7,4$ & -- & $168,168,24$ & $2,2,2$ & exact \\
$Sp_4(2)$ on polar points / Lagrangians & $15$ & -- & $720$ & $3$ & exact \\
$GL_3(2)$ on complete flags & $21$ & $|B|=8$ & $168$ & $6$ & coset \\
\midrule
$M_2(\F_2)\rtimes(GL_2(2)\times GL_2(2))$ & $16$ & $36$ & $576$ & $3$ & affine \\
$\operatorname{Alt}_3(2)\rtimes GL_3(2)$ & $8$ & $168$ & $1344$ & $2$ & affine \\
$\operatorname{Herm}_2(4/2)\rtimes GL_2(4)$ image & $16$ & $60$ & $960$ & $3$ & affine \\
$\mathcal Q(2,3)\rtimes GL_2(3)$ image & $27$ & $24$ & $648$ & $5$ & affine \\
$\F_2^4\rtimes O_4^+(2)$ & $16$ & $72$ & $1152$ & $3$ & affine polar \\
Folded and halved $4$-cubes & $8$ & $24$ & $192$ & $3$ & affine \\
\midrule
Shrikhande full vertex action & $16$ & $12$ & $192$ & $4$ & graph \\
$4\times4$ rook graph & $16$ & -- & $1152$ & $3$ & graph \\
Doob $D(1,1)$ chosen product action & $64$ & -- & $4608$ & $8$ & product \\
\bottomrule
\end{longtable}
\end{center}

For the supplemental affine and graph cases, the reported low-degree values
are:
\[
\begin{array}{c|cccccccccc}
\text{case}&\mathrm{Bil}&\mathrm{Alt}&\mathrm{Herm}&\mathcal Q&\mathrm{AffPol}
&\mathrm{Fold}&\tfrac12Q&\mathrm{Sh}&\mathrm{Rook}&D(1,1)\\ \hline
[m_{11}]&3&2&3&5&3&3&3&4&3&8
\end{array}
\]
In every entry, the matrix-trace average, cycle formula, and direct orbit count
agree.  The quadratic example also illustrates why $[m_{(2)}]$ need not equal
$[m_{(1,1)}]$: its values are $4$ and $5$, respectively, because forgetting
the two colors identifies a difference with its negative.

\resultbox{\textbf{Verification outcome.}  The suite completes 99 exact
coefficient comparisons.  Every displayed coefficient agrees across all
implemented routes; every family-specific power-fixed-point formula
reconstructs the literal cycle counts; the $GL_3(2)/B$ example verifies the
centralizer/class-intersection formula element by element; and the Doob
$D(1,1)$ example verifies the product fixed-point identity for powers
$1\leq r\leq6$.}

\subsection{Scope and cautions}

\begin{itemize}
\item The universal theorem is valid for every finite action.  An association
scheme is useful structure for calculating $F_r(g)$; it is not an additional
hypothesis needed for the theorem.
\item The automorphism group must be specified.  Different subgroups of
$\Aut(\Omega)$ generally give different averages and different orbit counts.
\item Parameters such as an intersection array or strongly regular tuple do
not determine the automorphism action, as the rook/Shrikhande comparison
shows.
\item Semilinear, complement, or exceptional automorphisms are incorporated
by averaging over their extra cosets.  Omitting such a coset computes the
series for the smaller stated group, not for the full automorphism group.
\item Exact finite formulas, stable formulas, and coefficientwise limits are
distinct claims.  This note states a stable range only where bounded support
and extension give one directly.
\end{itemize}
